% =============================================================
%  Tokens on a Timeline: Tangible Correction of Automatic
%  Speaker Diarization for Low-Resource Speech Annotation
%
%  Target journal: Behaviour & Information Technology
%                  (Taylor & Francis, Q1)
%  Style: APA author-year
%  Notes for production:
%    - Replace \documentclass{article} with the official
%      Taylor & Francis class \documentclass{tandfarticle}
%      when the journal's cls is available locally.
%    - Move authorship block to the journal's preferred
%      \author{...} \affiliation{...} macros for camera-ready.
% =============================================================

\documentclass[10pt,a4paper,twocolumn]{article}

% --- Packages ----------------------------------------------------
\usepackage[utf8]{inputenc}
\usepackage[T1]{fontenc}
\usepackage[margin=0.75in]{geometry}
\usepackage{times}
\usepackage{microtype}
\usepackage{setspace}
\usepackage{graphicx}
\usepackage{booktabs}
\usepackage{hyperref}
\usepackage{url}
\usepackage{xcolor}
\usepackage[round,authoryear]{natbib}
\bibliographystyle{apalike}

% --- Title and authors -------------------------------------------
\title{\textbf{Tokens on a Timeline: Tangible Correction of Automatic
       Speaker Diarization for Low-Resource Speech Annotation}}

\author{
  Arnab Satyam Roy \\
  Department of Computer Science \\
  American International University--Bangladesh (AIUB) \\
  Dhaka, Bangladesh \\
  \href{mailto:arnabsroy9@gmail.com}{arnabsroy9@gmail.com}
  \and
  Ahsanul Haque \\
  Department of Computer Science \\
  American International University--Bangladesh (AIUB) \\
  Dhaka, Bangladesh \\
  \href{mailto:ahsanulhaque1231@outlook.com}{ahsanulhaque1231@outlook.com}
  \and
  Abdus Salam \\
  Department of Computer Science \\
  American International University--Bangladesh (AIUB) \\
  Dhaka, Bangladesh \\
  \href{mailto:abdus.salam@aiub.edu}{abdus.salam@aiub.edu}
}

\date{\today}

\begin{document}

\twocolumn[\begin{@twocolumnfalse}
\maketitle
\end{@twocolumnfalse}]
\onehalfspacing

% =============================================================
%  1. Introduction
% =============================================================
\section{Introduction}\label{sec:intro}

Long-form speech understanding, comprising both transcription and
speaker attribution, underpins applications including media
indexing, meeting summarization, accessibility tooling, and legal
and forensic audio analysis. Current systems remain optimized for
high-resource languages, and their performance on conversational,
multi-speaker, low-resource audio is markedly weaker. Recent work
on Bangla long-form speech indicates that further gains depend on more annotated
diarization data: the diarization stage of an end-to-end pipeline is
limited by a small number of annotated recordings \citep{bhuiyan2026banglawsd}, and out-of-the-box
diarization on Bangla conversational content reaches a
diarization error rate (DER) of around 40\% on the only public
benchmark \citep{tabib2026bengaliloop}.

This constraint points to an opportunity located in the interaction
layer rather than the model. Faster and more reliable human
correction tools yield more corrected annotation, which in turn
supports better models for under-resourced languages. The
contribution proposed here lies in that interaction layer. The
diarization model is frozen and treated as a fixed backend, and no
modeling contribution is claimed. Bangla provides the motivation
and the deployment setting, namely annotation scarcity, while the
interaction question concerns the general problem of correcting
the temporal output of a learned recognizer.

We describe a tangible user interface in which a physical token
represents each speaker and is placed on a printed timeline to
correct the output of a frozen pre-trained diarization model. The
contribution is an interaction technique, the empirical hypothesis
that justifies it, and the study design that will test the
hypothesis. Section~\ref{sec:lit} reviews related work in four
threads: the automatic diarization pipeline and the human-facing
correction bottleneck, tangible interaction for time-based
information, empirical comparisons of tangible and graphical
modalities, and the low-resource Bangla speech setting.
Section~\ref{sec:gap} states the gap addressed at the intersection
of these threads. Section~\ref{sec:obj} lists the objectives and
hypotheses. Section~\ref{sec:method} specifies the proposed system
and study design in full.

% =============================================================
%  2. Literature Review
% =============================================================
\section{Literature Review}\label{sec:lit}

\subsection{Automatic speaker diarization and the correction
  bottleneck}
\label{sec:lit-diarization}

Speaker diarization is the task of partitioning an audio stream
into homogeneous segments according to speaker identity
\citep{bredin2020pyannote}. Contemporary systems follow a
segmentation--embedding--clustering pipeline in which a
sliding-window segmenter produces frame-level speaker
activities, a per-window embedding model summarises each
window into a fixed-dimensional vector, and an agglomerative
or spectral clustering groups windows into speaker tracks. The
pyannote.audio toolkit is a widely used open-source
reference implementation, with pre-trained models for voice
activity detection, speaker change detection, overlapped speech
detection, and speaker embedding \citep{bredin2020pyannote}.
Recent work on the loss formulation shows that switching from
the standard multi-label encoding to a powerset multi-class
encoding, in which each overlap pattern is a separate class,
removes the fragile detection threshold and yields consistent
gains on a wide range of benchmarks \citep{plaquet2023powerset}.
Long-form recognition is increasingly delivered by weakly
supervised encoder--decoder models such as Whisper
\citep{radford2023whisper}, which has become a common backbone
for Bangla long-form ASR \citep{bhuiyan2026banglawsd}. The
end-to-end Bangla pipelines that combine Whisper and pyannote
report substantial improvements over the unadapted baselines,
but the same line of work points to annotated data, not
modelling, as the limiting factor for further gains.

Recent work has explored four routes to making the correction step
cheaper or faster. The first is post-hoc, model-driven correction of
the diarization output. \citet{he2026inmeeting} describe an
LLM-assisted workflow in which a user provides brief corrective
feedback in natural language, an LLM interprets which span to
re-assign, and online speaker enrollments are added to prevent
future attribution errors. A complementary direction uses a
fine-tuned LLM as a speaker error corrector at the transcript
level, taking the initial diarization hypothesis and the ASR
transcript as input and producing a corrected label sequence
\citep{kumar2025seal, efstathiadis2024llmdiar}. The latter
line demonstrates that LLM-based correction can markedly improve
diarization accuracy on a held-out benchmark, but the gains are
constrained to transcripts produced by the same ASR tool that
generated the fine-tuning data. Related work reformulates the
correction as an end-to-end neural post-processing step that
exploits the interaction between the initial speaker activity
and the acoustic features \citep{han2024diacorrect}. All of
these systems automate the correction step and so remove the
human from the loop, but they require supervised training data
that is itself scarce for low-resource languages, and they do
not address the underlying question of how a human annotator
can correct a model output most effectively.

The second direction is human-in-the-loop active correction. In
the diarization setting, this is exemplified by incremental
collection-level systems in which a human expert answers simple
yes/no questions about whether two segments were spoken by the
same speaker, and the system updates the clustering hypothesis
accordingly \citep{prokopalo2022active}. The questions are
selected by an algorithm that aims to minimise human effort
while maximising improvement in the diarization error rate, and
a stopping criterion caps the total interaction cost. The same
line of work generalises the framework to cross-show
diarization through a shared speaker database.

The third direction is the development of annotation tools that
instrument the cost of correction rather than only the resulting
accuracy. \citet{broux2018cassid} introduced the Human--Computer
Interaction Quantity (HCIQ), a weighted sum of correction-action
counts that captures the labour of a human correction session
without relying on the diarization error rate alone, and they
validated it on the REPERE French broadcast-news corpus. The
ALLIES corpus, a meta-corpus of French TV and radio recordings
that aggregates several earlier evaluation campaigns, was
explicitly designed for human-assisted diarization and uses a
similar cost-aware protocol \citep{tahon2024allies}. More
recently, \citet{yamaguchi2026voxmap} presented
voxmap-studio, a React-based open-source diarization annotation
tool that records edit-operation counts and active time as a
first-class output, and used a three-condition preliminary
study on the AMI corpus to compare forms of within-modality
assistance. A separate semi-automatic toolkit, AnnoTheia,
addresses audio-visual speech annotation in low-resource
languages by combining active-speaker detection,
transcription, and human-in-the-loop verification
\citep{acosta2024annotheia}. The present article is closest in
spirit to this line of work: it treats annotation cost and
accuracy as joint outcomes, but it varies the interaction
modality rather than the form of within-modality assistance.

The fourth direction is the model-level refinement that sits
behind any human-facing tool. This includes both loss-function
improvements \citep{plaquet2023powerset} and end-to-end Bangla
pipelines \citep{bhuiyan2026banglawsd}. The present article
treats the backend as fixed and isolates the interaction
question.

\subsection{Tangible interaction for temporal media}
\label{sec:lit-temporal-tui}

A separate thread of work embodies time physically. The reacTable
\citep{jorda2007reactable} demonstrated that tabletop tangible
surfaces can host rich temporal interaction through object
manipulation. ChronoTape \citep{bennett2012chronotape} presents a
tangible timeline for personal and family history, in which a
paper tape is annotated through a reader device. Tquencer
\citep{kaltenbrunner2018tquencer} provides tangible control for
time-based media sequencing through passive tokens and overlays.
The Tangible Video Editor \citep{zigelbaum2007tve} supports
collaborative video editing through a set of active physical
tokens that represent clips and transitions.
\citet{mahieux2022sablier} use the hourglass as a tangible
metaphor for navigating through space and time in cultural
heritage applications, with a co-design study that demonstrates
acceptability among museum professionals. These systems author
or arrange temporal media; none addresses the correction of
output produced by a learned temporal recognizer.

Theoretical grounding for token-based temporal interfaces comes
from the token+constraint framework, in which physical tokens
represent digital information and physical constraints are mapped
to digital operations \citep{ullmer2005tokenconstraint}. The
related Tangible Query Interfaces apply the same idea to database
query construction \citep{ullmer2003tangiblequery}. Both build on
Ishii and Ullmer's broader ``Tangible Bits'' vision of bridging
bits and atoms through graspable physical artefacts
\citep{ishii1997tangiblebits}. A complementary thread
investigates the design of passive tokens for multi-touch
surfaces: \citet{morales2016touchtokens} show that the geometry
of a token determines the grasp users adopt, and that the
resulting touch pattern on the surface can be recognised with
high accuracy, providing a low-cost way to combine tangible and
gestural input without instrumenting the surface itself. The
present study inherits the design principle that the token
itself, not the visual channel, carries the identity anchor,
and applies it to a temporal-correction task rather than to
multi-touch gesture disambiguation.

None of the systems reviewed above treats the token as a
per-instance identity anchor for the purpose of correcting a
recogniser's temporal output, which is the operation the
present article investigates. The closest relative is the use
of physical objects as identity anchors in tangible database
queries, but the queries there are constructed from scratch
rather than corrected against an automatically generated
hypothesis, and the temporal dimension is not central.

\subsection{Empirical comparisons of tangible and graphical
  interaction}
\label{sec:lit-tui-vs-gui}

A small but methodologically relevant body of work directly
compares tangible and graphical user interfaces on a common
task. The earliest formal comparison in the HCI literature is
\citet{patten2000spatcmp}, who contrasted a token-based TUI
with a graphical icon-based GUI on a location-recall task and
reported an advantage for the tangible condition. The TUI
subjects performed better at recall and were also more likely
to use absolute spatial reference frames to organise the
blocks, a strategy that correlated with better recall. The
work is a methodological model for the present study: a
within-subjects comparison on a single task, with both
conditions matched in capability and presentation.

A second strand measures the cognitive cost of tangible
interaction. \citet{chandrasekera2015cogload} compared a
desktop augmented-reality TUI with a desktop virtual-reality
GUI on a creative design task and found significantly lower
self-reported cognitive load in the TUI condition
($p = 0.045$), measured with the NASA Task Load Index
(NASA-TLX). The interpretation offered is that the tangible
interface supports epistemic action, the use of physical
manipulation to off-load cognitive work, more effectively than
the graphical interface. The present article draws on this
methodological lineage: it uses NASA-TLX as a secondary
outcome and tests a comparable primary question (correction
accuracy) on a different task (diarization correction).

A third strand compares tangible and graphical interaction in
the context of the token+constraint framework
\citep{ullmer2003tangiblequery, ullmer2005tokenconstraint}.
Ullmer et~al.'s preliminary user study of a tangible database
query interface against a comparable graphical interface
reported a measurable advantage for the tangible condition on
the specific sub-task where physical rotation mapped to
parameter value changes; the advantage disappeared for
sub-tasks that did not exploit a physical affordance unique to
the tangible condition. The present study follows the same
logic of designing the comparison around the affordance unique
to each modality, and the double-dissociation prediction in
Section~\ref{sec:obj} is an explicit generalisation of the
affordance-specific-advantage pattern.

\subsection{Low-resource Bangla speech as the deployment setting}
\label{sec:lit-bangla}

The deployment context for the proposed work is Bangla, a
low-resource language in the long-form speech regime. Public
Bangla corpora remain dominated by short, single-speaker
utterances collected through crowdsourcing, and the largest
publicly available corpus of this kind is the 229-hour Google
Bangla speech corpus \citep{kjartansson2018crowdsourced}. The Bengali-Loop benchmark
\citep{tabib2026bengaliloop} was introduced specifically to
address the long-form gap and consists of two sub-corpora: a
long-form ASR corpus of 191 recordings totalling 158.6 hours
from eleven YouTube channels, with a reproducible
subtitle-extraction and human-in-the-loop verification pipeline,
and a speaker diarization corpus of 24 recordings totalling 22
hours with fully manual speaker-turn labels in CSV format. The
benchmark reports that an out-of-the-box pyannote.audio pipeline
reaches a diarization error rate of 40.08\% on the diarization
test set, and that a BengaliAI fine-tuned Whisper model reaches
a word error rate of 34.07\% on the ASR test set. These numbers
anchor the gap between unadapted performance and the level
required for downstream applications, and they are consistent
with the claim that further gains on low-resource languages
depend primarily on labelled data rather than on modelling
innovation.

The DL Sprint 4.0 Kaggle competition, run on a subset of the
Bengali-Loop data, attracted multiple community submissions
that fine-tuned the pyannote segmentation backbone and the
tugstugi Whisper model on the competition training set. One
submission \citep{bhuiyan2026banglawsd} reports a DER of
23.92\% on the diarization track and a WER of 24.41\% on the
ASR track, both well below the unadapted baselines; other
submissions report even lower numbers, with public-leaderboard
DER around 20\% and public WER around 16\%
\citep{tanvir2026shobdosetu}. The competition reports also produced three
pieces of empirical evidence that are directly relevant to the
present work: first, that fine-tuning the segmentation component
of pyannote.audio is more impactful than fine-tuning the
embedding model \citep{chowdhury2026robustbn}; second, that vocal source
separation combined with silence-aware chunking reduces
word-boundary errors by preserving sentence-internal structure
\citep{chowdhury2026robustbn}; and third, that targeted data augmentation
for muffled-zone robustness contributes measurable gains
\citep{tanvir2026shobdosetu}. Each of these findings depends on having
access to labelled long-form Bangla audio, and the labelled
data was produced in part by manual annotation. The competition
submissions themselves include a description of the data
pipeline used to construct the training set, and the manual
verification of the automatically extracted transcripts and
the manual annotation of speaker turns is reported as a
substantial component of that pipeline. The present work is
complementary to these modelling efforts: by targeting the
human annotation step, it attacks the constraint that the
modelling efforts also identify, namely the scarcity of
high-quality annotated data.

% =============================================================
%  3. Research Gap
% =============================================================
\section{Research Gap}\label{sec:gap}

The four threads reviewed in Section~\ref{sec:lit} have evolved
largely independently. Screen-based diarization correction is
mature, with both commercial and research tools that instrument
correction cost \citep{broux2018cassid, tahon2024allies,
yamaguchi2026voxmap}; tangible interfaces for time concern
sequencing and authoring of temporal media
\citep{jorda2007reactable, bennett2012chronotape,
kaltenbrunner2018tquencer, mahieux2022sablier}; and human-in-the-loop
annotation has been shown to be effective when the interaction is
graphical \citep{prokopalo2022active}. The token+constraint
framework provides a vocabulary for tangible interaction with
digital information \citep{ullmer2005tokenconstraint,
ullmer2003tangiblequery, ishii1997tangiblebits}, and empirical
work has demonstrated measurable differences between tangible and
graphical interaction on memory and load outcomes
\citep{patten2000spatcmp, chandrasekera2015cogload}.

The gap addressed in this article lies at the intersection of
these areas. The specific combination of (i) passive token-based
correction, (ii) of the temporal output produced by a frozen
recognizer, (iii) deployed as an annotation accelerator for a
low-resource language, has not been explored. No existing study
asks whether the choice of interaction modality governs which
class of recognizer error a human annotator catches.
The closest contemporary neighbour
\citep{yamaguchi2026voxmap} varies graphical assistance while
holding modality constant. The present study holds the recognizer
and the annotation operations constant and varies the modality.

% =============================================================
%  4. Objectives
% =============================================================
\section{Objectives}\label{sec:obj}

The central question is whether the interaction modality used to
correct automatic speaker diarization changes which class of error
a human annotator catches and repairs, and at what cost in time
and effort. The study tests three hypotheses.

\textbf{H1.} Annotators using the tangible interface correct
speaker-confusion errors, in which speech is attributed to the
wrong speaker, with higher accuracy and lower correction time than
with the graphical interface, because per-speaker tokens
externalize identity tracking.

\textbf{H2.} Annotators using the graphical interface correct
temporal boundary errors, that is the offsets of segment start
and end points, with greater precision than with the tangible
interface, because, at a temporal scale held identical across
conditions, the graphical interface affords finer temporal control
through continuous pointer positioning on the waveform, whereas the
tangible interface is limited by the positional resolution of the
marker tracker.

\textbf{H3.} Perceived workload and stated preference differ by
modality. These are treated as secondary outcomes.

The primary prediction is the interaction between modality and
error type, the double dissociation, which is more robust to
novelty and engagement effects than any single main effect. A
confirmatory result on the interaction, with H1 and H2 pointing in
opposite directions, would establish that interaction modality
governs which class of recognizer error a human annotator catches,
a principle that transfers beyond diarization to other temporal
labeling tasks. Detecting a crossover interaction of moderate
size requires more statistical power than detecting a single
main effect; the sample size in Section~\ref{sec:method-design}
is therefore set with the interaction term as the primary
inference target, and a sensitivity analysis is pre-registered
for the smallest effect of the predicted direction that the
study can reliably detect.

% =============================================================
%  5. Proposed Methodology
% =============================================================
\section{Proposed Methodology}\label{sec:method}

\subsection{Tangible interface}
\label{sec:method-tangible}

The interface uses one passive token per speaker. Each token is a
printed ArUco fiducial marker fixed to a small rigid object such
as a wooden craft cube or a bottle cap, which requires no
specialized fabrication. The camera reads only marker identity
and position. Tokens are placed on a printed paper timeline
divided into horizontal lanes. Placing a speaker's token in a
segment assigns that segment to the speaker, and moving the token
to another segment reassigns it. Boundaries are edited with a
separate boundary-handle token that snaps to the nearest segment
edge in the lane on which it rests, and sliding the handle moves
that edge. The two operations are carried by different physical
tokens and are told apart by marker identity, so a reassignment
and a boundary edit never share a gesture. The handle gives the
tangible condition a genuine if coarser boundary-editing
capability, so that the boundary comparison is not decided by
capability alone.

Sensing uses a consumer USB webcam of at least 1080p resolution,
mounted overhead on a copy stand and feeding directly into
OpenCV. A local fixed-framerate feed avoids the latency and
connection fragility of a phone-camera arrangement. Autofocus is
disabled and exposure and white balance are fixed through the
camera controls to prevent tracking loss. Because webcam
resolution determines marker positional resolution, and therefore
the smallest boundary adjustment the tangible condition can
express, the camera is fixed and characterized before data
collection. Marker detection is restricted to identity and
position, with no rotation or multi-token grammar, which keeps the
vision pipeline small and inexpensive to run.

During the tangible condition the screen presents the current
waveform, segment boundaries, speaker coloring, and playback
position, but accepts no input. The mouse and keyboard are
removed from the workspace so that the screen cannot serve as an
alternative interface, and all correction is performed through
the tokens. Feedback is therefore visual rather than haptic,
consistent with established tabletop tangible interfaces such as
the reacTable \citep{jorda2007reactable}.

\subsection{Graphical baseline}
\label{sec:method-graphical}

The graphical baseline is built by wrapping the open-source
waveform library wavesurfer.js, together with its regions plugin,
in a minimal interface. This reuses the most demanding parts of
the frontend, including waveform rendering, zooming, playback
synchronization, and draggable region handles,
while keeping the comparison under experimental control. A full
annotation platform is deliberately not adopted, because its
accumulated features would introduce differences unrelated to
interaction modality and reducing it to parity would often exceed
the effort of a focused implementation.

The baseline is governed by three constraints. It exposes exactly
the operations available in the tangible condition, namely
reassigning a segment to a speaker by selecting its region and
adjusting a boundary by dragging a region edge, and no others. It
omits production features that have no tangible counterpart,
including keyboard shortcuts, multi-track layouts, snapping,
autosave, and free zooming of the waveform, whose temporal scale
is instead fixed and matched across conditions
(Section~\ref{sec:method-control}). It remains a competent editor, since a weakened
baseline would invalidate any observed advantage for the tangible
interface.

\subsection{Diarization backend}
\label{sec:method-backend}

A pre-trained Bangla diarization pipeline, comprising pyannote
segmentation, speaker embedding, and clustering
\citep{bredin2020pyannote, plaquet2023powerset}, produces
diarization output before any session. No model runs during
interaction, so there is no computational load at interaction time
and laptop-grade hardware is sufficient. The model is not
modified and is not part of the contribution.

\subsection{Experimental control across conditions}
\label{sec:method-control}

Both conditions are implemented as a single application with a
shared state model and a shared renderer, differing only in the
source of input. The on-screen representation is produced by the
same rendering code in both conditions, which removes visual
presentation as a source of difference. In the graphical
condition the input source is mouse interaction with the rendered
regions. In the tangible condition the input source is the marker
tracker, which emits the same logical operations of speaker
reassignment and boundary movement while the rendered regions
remain non-interactive. Sharing the waveform view across
conditions is also what makes the boundary comparison
interpretable. The waveform is drawn at a single fixed temporal scale that shows
the whole clip, and interactive zooming is disabled in both
conditions, so neither modality has finer visual access to the
signal than the other. Because both conditions act against the
same fixed visual reference, the boundary measurement reflects the
precision of the manipulation channel, boundary-handle sliding
against pointer dragging, rather than differences in visual
access. A zoom-assisted graphical editor would carry a resolution
advantage that belongs to the tool rather than to the interaction
modality, and fixing the scale removes that confound while
preserving the finer positioning the pointer affords at that
scale. The fixed scale and the predicted direction of the boundary
effect are set before data collection.

Playback is held constant across conditions through a single
physical control used for play, pause, and seeking in both, so
that playback does not introduce an asymmetry unrelated to the
interaction modality.

Two consequences of strict information parity are noted. First,
because speaker identity is shown on screen in both conditions, a
confirmed H1 is attributable to the manipulation of identity,
that is grasping and placing the correct token without visual
search of a label, rather than to a stronger form of external
representation. This narrows the mechanism while strengthening
internal validity. Second, with a vertical monitor the locus of
action on the table is spatially separated from the locus of
feedback on the screen, whereas the graphical condition
co-locates them. This separation is a property of the tangible
setup as instantiated and is reflected in the scope of the
claims, which concern token interaction on a table with
screen-based feedback.

\subsection{Design and participants}
\label{sec:method-design}

The study uses a within-subjects design with two conditions, the
tangible interface and the graphical baseline, and counterbalanced
order to control for learning effects. Each participant corrects
a set of diarized audio clips under each condition. Detecting an
interaction requires more statistical power than detecting a
single main effect, so the study targets between 16 and 24
participants. The faculty collaborator conducts recruitment and
runs the sessions.

\subsection{Stimulus construction}
\label{sec:method-stimulus}

The error distribution of the frozen model is first characterized
on held-out Bangla audio, measuring the relative frequency of
speaker-confusion errors and boundary errors. Clips are then
constructed by injecting a controlled and balanced set of both
error types onto ground-truth diarization, at rates consistent
with the model's actual behavior, so that the exact errors are
known in advance and matched across conditions and participants.
Injected boundary offsets are sampled from the model's measured
offset distribution, with offsets smaller than the characterized
tangible resolution floor excluded, so that every injected
boundary error is large enough to be corrected in either condition
and the comparison is not confounded by errors that only one
modality can physically express.
Total error count and clip difficulty are held constant across
trials, and clip sets are counterbalanced across conditions. The
number of speakers is held constant to reduce variance.

\subsection{Measures}
\label{sec:method-measures}

The primary measures are correction accuracy for each error type,
defined as the proportion of injected errors of that type that
are correctly repaired, and correction time for each error type.
Boundary precision is measured as the absolute timestamp error
remaining after correction. Secondary measures are perceived
workload, assessed with the NASA Task Load Index
\citep{hart1988nasatlx}, following the precedent of
\citet{chandrasekera2015cogload}, and stated preference.

\subsection{Analysis}
\label{sec:method-analysis}

The primary analysis is a two-factor repeated-measures model with
modality and error type as factors, applied to accuracy and to
correction time. The interaction term tests the predicted
dissociation. Boundary precision is analyzed separately.
Secondary measures are reported with appropriate non-parametric
tests.

% =============================================================
%  References
% =============================================================
\bibliography{ref}

\end{document}
